\section{Collision Probabilities of Centera CA}\label{sec:5}

Centera has three naming schemes, referred to as \textbf{M}, \textbf{GM}
and \textbf{M++.} The \textbf{M} naming scheme returns a 128-bit CA that
is simply the MD5 hash of the associated content. The \textbf{GM} naming
scheme returns a 256-bit CA that is comprised of an MD5 hash of the
associated content, a random bit-string, timestamp, counter and header
bits. The \textbf{M++} naming scheme returns a 256-bit CA that is
comprised of the MD5 and a truncated SHA-256 hash of the associated
content. Detail for \textbf{M} is provided in \S\ref{sec:5.2};
\textbf{GM} detail is provided in \S\ref{sec:5.4}; \textbf{M++}
is treated in \S\ref{sec:5.5}. For now the salient point is that
the various naming schemes return CAs of different bit size (128 bits
for \textbf{M} and 256 bits for \textbf{GM} and \textbf{M++}), and hence
have differing collision probabilities and preimage resistance.

\subsection{Birthday Paradox}\label{sec:5.1}\label{birthday-paradox}

Question: Ignoring leap years, how many people need to be present in a
room for there to be \(\approx50\%\) chance that someone in the room
shares (i.e. collides with) your birthday? Answer: 253 (as shown in
\eqref{eq:5.1} below)\footnote{This maps to a preimage or
  second-preimage hash collision.}.

\begin{equation}\label{eq:5.1}
1-\left(\frac{364}{365}\right)^{253}\approx0.50
\end{equation}

However, if you change the question to include one key difference, the
answer changes dramatically.

The modified question is: how many people need to be present in a room
for \emph{any two} people to have \(\approx50\%\) chance of sharing the
same birthday? Answer: 23. This number is surprisingly low, which leads
to this phenomenon being given appellations such as ``birthday
surprise'' or ``birthday paradox.'' The surprisingly small result for
the second question becomes clear when you consider that with 23 people
in the room there are 253 \emph{pairs} of people in the room; i.e.,
\(23\cdot22/2=253.\) Of course, for this result to hold the people in
the room must have an even (random) distribution of birthdays; clearly
this result would not stand if the room was full of twins.

We now add some rigor and generalization to the birthday paradox.

\textbf{Theorem 3:} Let \(P_{\mathrm{collision}}(N,q)\) denote the
probability to have at least one collision when we throw \emph{q
\textgreater{}} 1 balls at random into \emph{N \textgreater{} q}
buckets, giving us equation \eqref{eq:5.2} below.

\begin{equation}\label{eq:5.2}
P_{\mathrm{collision}}(N,q)\le0.5\frac{q(q-1)}{N}
\end{equation}

In this paper we only present an upper bound to
\(P_{\mathrm{collision}}(N,q)\) since that is relevant for our collision
probability calculations. For a more exhaustive proof, including a lower
bound, readers are referred to \cite{ref26}.

\textbf{Proof:} Let \(C_i\) be the event that the \(i\) -th ball is
thrown in a non-empty bucket, i.e.~it collides with one of the previous
balls. Then \(P[C_i]\) is at most \((i-1)/N\) since when the \emph{i}-th
ball is thrown, there are at most \(i-1\) different occupied buckets and
the \(i\) -th ball is equally likely to land in any of them. Now we
have,

\begin{equation}\label{eq:5.3}
\begin{aligned}P_{\mathrm{collision}}(N,q)&=P[C_1\lor C_2\lor\cdots\lor C_q]\\&\le P[C_1]+P[C_2]+\cdots+P[C_q]\\&\le\frac0N+\frac1N+\frac2N+\cdots+\frac{q-1}{N}\\&\le\frac{q(q-1)}{2N}.\end{aligned}
\end{equation}

From equation \eqref{eq:5.3} we see the quadratic relation
between \emph{q} and \(P_{\mathrm{collision}}(N,q).\) This is the core
of the birthday paradox. You only need to throw about \(\sqrt{N}\) balls
into \emph{N} buckets to achieve a reasonable probability of ending up
with two (or more) balls in the same bucket.

\subsection{Centera with the M naming
scheme}\label{sec:5.2}\label{centera-with-the-m-naming-scheme}

If a large number of messages are hashed using a given hash function,
collisions will occur ``by chance'' after a very large number of
messages have been hashed. Extensive statistical tests have been
performed in the past on MD5 (for example in the RIPE project
\cite{ref19}), and no deviations from random behavior have
been detected. This implies that modeling MD5 in this respect as a
random function will result in a very good approximation of the
collision probability.

MD5 hashes data into 128-bit strings. Relating to the proof in
\eqref{eq:5.3} we can state that \(N=2^{128}.\) Therefore, once
a threshold of \(q=2^{64}\) files is reached, collisions will be found
very quickly. More precisely, if \(q=\alpha\cdot2^{64}\) files are
hashed, then we can derive from formula \eqref{eq:5.2} that the
probability that one or more collisions occur is upper bounded by
\eqref{eq:5.4}.

\begin{equation}\label{eq:5.4}
\frac{\alpha\cdot2^{64}\cdot(\alpha\cdot2^{64}-1)}{2\cdot2^{128}}\approx\frac{\alpha^2}{2}
\end{equation}

Note that Centera typically operates with very small values of
\(\alpha\) (e.g., .000000001) since the total file count in the system
is substantially less than \(2^{64}\) files (18 billion billion files).

Table~\ref{tab:1} indicates the probability of a CA collision when
using the \textbf{M} naming scheme for various fill rates.

\begin{table}[htbp]
\centering\small
\begin{tabularx}{\linewidth}{@{}r*{3}{>{\raggedright\arraybackslash}X}@{}}
\toprule
Number of objects & Capacity with avg file size 10 bytes & Capacity with avg file size 1 kB & Approximate M collision probability\\
\midrule
$10^6$ & 10 Megabyte & 1 Gigabyte & $1.47\times10^{-27}$\\
$10^7$ & 100 Megabyte & 10 Gigabyte & $1.47\times10^{-25}$\\
$10^8$ & 1 Gigabyte & 100 Gigabyte & $1.47\times10^{-23}$\\
$10^9$ & 10 Gigabyte & 1 Terabyte & $1.47\times10^{-21}$\\
$10^{10}$ & 100 Gigabyte & 10 Terabyte & $1.47\times10^{-19}$\\
$10^{11}$ & 1 Terabyte & 100 Terabyte & $1.47\times10^{-17}$\\
$10^{12}$ & 10 Terabyte & 1 Petabyte & $1.47\times10^{-15}$\\
$10^{13}$ & 100 Terabyte & 10 Petabyte & $1.47\times10^{-13}$\\
$10^{14}$ & 1 Petabyte & 100 Petabyte & $1.47\times10^{-11}$\\
$10^{15}$ & 10 Petabyte & 1 Exabyte & $1.47\times10^{-9}$\\
\bottomrule
\end{tabularx}
\caption{Accidental collision probability for M; decimal logical payload capacities}\label{tab:1}
\end{table}


Such numbers are sufficiently large that it's difficult to get a sense
of magnitude. Using time as the dimension, assume a hypothetical CAS
application capable of storing 10,000 files per second to a Centera
cluster. Using an average file size of 1K, writing \(10^{15}\) such
files would take up an exabyte of protected storage. However, it would
take several millennia to fill up the cluster at this rate. As a point
of comparison, after 1,000 years of continuous operation the probability
of a CA collision using the M naming scheme is less than 1 in
1,000,000,000.

In the case of a malicious user seeking to forge two arbitrary files
which return an identical \textbf{M} content address (this is referred
to as ``collision resistance''), then the user must attempt to apply the
work described in \cite{ref30}\cite{ref31} to
Centera.

In the case of a malicious user seeking to forge a file \(X'\) which
gets an identical \textbf{M} content address as a given file X (this is
referred to as ``preimage resistance''), one can attempt to apply the
ideas described in \cite{ref11}. Knowing that the maximal file
size in Centera is 100Mb, we have less than \(2^{21}\) 512-bit blocks in
one file. We can conclude that on average about \(2^{128-21+1}=2^{108}\)
MD5 evaluations are required to forge a preimage of a given file X.
Although this effort can happen offline, the number is sufficiently
large to make such attacks impractical.

\subsection{Random Generators}\label{sec:5.3}\label{random-generators}

In order to move towards the stronger GM naming scheme, we need the
following definitions from the HAC \cite{ref2}.

\emph{Definition 6}: A \textbf{pseudorandom bit generator} (PRBG) is a
deterministic algorithm which, given a truly random binary sequence of
length \emph{k}, outputs a binary sequence of length \(l\ge k\) which
``appears'' to be random. The input to the PRBG is called the
\emph{seed}, while the output of the PRBG is called a \emph{pseudorandom
bit sequence}.

\emph{Definition 7}: A pseudorandom bit generator is said to pass the
\textbf{next-bit test} if there is no polynomial-time algorithm which,
on input of the first \emph{l} bits of an output sequence \emph{s}, can
predict the \emph{(l+1)}\textsuperscript{\emph{st}} bit of \emph{s} with
probability significantly greater than \(\frac12\).

\emph{Definition 8}: A PRBG that passes the next-bit test is called a
cryptographically secure pseudorandom bit generator (CSPRBG).

\subsection{Centera with the GM naming
scheme}\label{sec:5.4}\label{centera-with-the-gm-naming-scheme}

Although the probabilities of having collisions with the M naming scheme
are minimal, a second naming scheme, called GM, was introduced which
includes additional random bits in the resultant CA. This additional
randomness has the effect of further reducing the probability of
collisions.

The Centera GM CA contains the following components:

\begin{itemize}
\tightlist
\item
  \textbf{M:} a 128 bit MD5 hash based on the file content;
\item
  \textbf{G:} a 70 bit random string generated by a cryptographically
  secure pseudorandom bit generator;
\item
  \textbf{T:} a 35 bit timestamp, with second resolution (actually, one
  unit of \textbf{T} is 1024 milliseconds);
\item
  \textbf{C:} a 10 bit counter; and
\item
  \textbf{H:} 13 bits used as ``header'' codes.
\end{itemize}

\subsubsection{Timestamp and counter}\label{sec:5.4.1}\label{timestamp-and-counter}

The timestamp \textbf{T} significantly reduces the probability of
collisions by effectively partitioning the entire Centera system into
sets of files that have been written within about one second of each
other. Incoming files are time stamped by the access node through which
the file enters the cluster the first time. The timestamp \textbf{T}
overflows after a little more than 1,000 years.

Suppose that the peak write speed to a Centera is \emph{S} files per
second per access node. Given \emph{A} access nodes in a Centera, we
conclude that the upper bound for the number of files having identical
timestamp \textbf{T} is equal to \(A\cdot S.\)

Every access node in Centera maintains a 10-bit counter \textbf{C} which
is increased for every new file stored on the system. The value of this
counter \textbf{C} is randomized across the access nodes and repeats
after each 1,024 files stored through the access node. An upper bound
for the number of files having both identical timestamp \textbf{T} and
identical counter \textbf{C} is \eqref{eq:5.5}.

\begin{equation}\label{eq:5.5}
A+\left\lfloor\frac{A\cdot S}{2^{10}}\right\rfloor\approx A+\frac{A\cdot S}{2^{10}}
\end{equation}

This formula holds for \emph{S} larger than \emph{A}, which is always
the case. For very large values of \emph{S}, the addition of \emph{A}
can be omitted as having minor impact on the result. Note that the
\emph{A} term is due to a quantization effect because it is not possible
to split documents into pieces for access node processing (e.g., when
1025 documents are assigned C values from the 1024 possibilities, at
least one value will be assigned to at least 2 documents).

\subsubsection{Applying the Birthday
Paradox}\label{sec:5.4.2}\label{applying-the-birthday-paradox}

In order for a collision to occur using the GM naming scheme, all
components \textbf{M.G.T.C.H} of the blob address must be identical for
two distinct files.

In \S\ref{sec:5.4.1} we showed that the impact of the timestamp
\textbf{T} and counter \textbf{C} effectively limits the size of the set
in which we can have collisions, to approximately \(A\cdot S/2^{10}\)
files (assuming that \emph{S} is much larger than \emph{A}).

The component \textbf{H} is non-random, hence it does not help to reduce
the probability of collisions.

The components \textbf{M} and \textbf{G} are fully random and add up to
128+70=198 bits (see \S\ref{sec:5.2} for the randomness of
\textbf{M} and \S\ref{sec:5.4} for the randomness of \textbf{G}).

Applying the birthday paradox to a set of \(A\cdot S/2^{10}\) balls
thrown into \(2^{198}\) buckets, we end up with the probability of a
collision with the GM naming scheme indicated in \eqref{eq:5.6}.

\begin{equation}\label{eq:5.6}
\frac{(A\cdot S)^2}{2^{20}\cdot2\cdot2^{198}}=\frac{(A\cdot S)^2}{2^{219}}
\end{equation}

The number of possible combinations (\textbf{T}, \textbf{C}) over a
large time interval is equal to the number of milliseconds in that time
interval.

If we run a Centera system for a period of \emph{Z} milliseconds, then
for large values of Z, the probability of experiencing a collision
during these \emph{Z} milliseconds is indicated in
\eqref{eq:5.7}.

\begin{equation}\label{eq:5.7}
\left(\frac{(A\cdot S)^2}{2^{219}}\right)Z
\end{equation}

\subsubsection{Collision Probability with
GM}\label{sec:5.4.3}\label{collision-probability-with-gm}

In a typical Centera setting, an upper bound for \emph{A} is 100
different access nodes, and an upper bound for \emph{S} is 10,000 small
files per second per access node. Running this experiment for 1,000
years, we get \(Z=3.15\cdot10^{13}\) milliseconds\emph{.} Feeding this
into equation \eqref{eq:5.7} we end up with the probability
described in the equation \eqref{eq:5.8} below.

\begin{equation}\label{eq:5.8}
\left(\frac{(100\times10{,}000)^2}{2^{219}}\right)Z\approx4\cdot10^{-41}
\end{equation}

In short, the probability of a collision in the GM naming scheme is
\(\sim4\cdot10^{-41}\) assuming an ingestion rate of 1,000,000
files/second over a period of 1,000 years. Note that this is an
extraordinarily high ingestion rate over a long period of time. In the
practical usage of Centera, the collision probability will be several
orders of magnitude smaller.

To help put this into perspective, the probability of a CA collision is
less likely than a bit error on disk (the non-recoverable Bit Error Rate
(BER) of a typical ATA disk is \textasciitilde1 for every \(10^{15}\)
bits read). As mentioned in §1, in addition to using the CA as an object
handle, Centera uses the CA as an MDC. Therefore, in the event of a bit
error on disk, Centera will detect this error via a background task that
periodically checks the validity of a file image on disk against its CA.
In the event a manipulation is detected, Centera will automatically
replace the bad file with a good copy, thus insuring data integrity over
the life of the archive.

In the case where a malicious user wishes to forge two arbitrary files
which result in an identical \textbf{GM} content address (CA), the user
needs to mount an online attack against a Centera system. One strategy
would be to start from a known \textbf{M} collision (as described in
\cite{ref30}) and attempt to write the colliding file to all
access nodes simultaneously within the same millisecond. In this way the
attacker would seek to get some control over the M.T.C.H components of
the \textbf{GM} content address (cf.~\S\ref{sec:5.4}). Within
that millisecond, the user would attempt to write to \emph{A} access
nodes in parallel and hope for a collision in the 70-bit G component.
Since \emph{A} is typically very small (less than 100), coupled with the
time it would take to perform the writes, it would take millions of
years before such an attack is successful.

For the same reasons, forging a second preimage in an attempt to
overwrite an existing \textbf{GM} file, is infeasible.

\subsection{Centera with the M++ naming
scheme}\label{sec:5.5}\label{centera-with-the-m-naming-scheme-1}

As discussed in \S\ref{sec:5.4}, the \textbf{GM} naming scheme as
implemented does not allow for single-instance store\footnote{More
  precisely, GM -- as presently implemented -- does not allow identical
  objects to have identical names. The \emph{function}, single-instance
  store, is anticipated to be activated for the GM naming scheme in a
  later release of CentraStar.}. For those applications that would like
to retain this ability, but use a stronger naming scheme than
\textbf{M}, a third naming scheme was created. This 256-bit naming
scheme, called \textbf{M++}, is constructed by the concatenation of 128
bits of MD5, 8 bits of formatting and 120 bits of SHA-256. The
calculation of object names in this naming scheme requires more CPU
cycles due to the introduction of the extra hash function evaluation.

Joux demonstrates that, contrary to Fact 9-27 in \cite{ref2},
the concatenation (or ``cascading'') of two iterated hash functions is
only as strong as the strongest of the two hash functions
\cite{ref1}. The following results apply to the \textbf{M++}
naming scheme.

When writing \textbf{M++} objects to a Centera, collisions become likely
after 2\textsuperscript{124} objects have been stored on the system
(cf.~the \textbf{M} naming scheme where this number is
2\textsuperscript{64}). This can be validated by looking at formula
\eqref{eq:5.2} and substituting \emph{q} with
2\textsuperscript{124} and \emph{N} with \(2^{120+128}=2^{248}\) to
reach a collision probability of \(\approx50\%\).

After 1,000 years of continuous operation, a hypothetical CAS
application capable of storing 10,000 files per second on a Centera
system would have stored \(3.15\cdot10^{14}\) files. This results in a
collision probability of 10\textsuperscript{-46} for the \textbf{M++}
naming scheme (cf.~10\textsuperscript{-9} for the \textbf{M} naming
scheme). This can be validated by looking at formula
\eqref{eq:5.2} and substituting \emph{q} with
\(3.15\cdot10^{14}\) and \emph{N} with \(2^{120+128}=2^{248}.\)

A malicious user wishing to induce a collision for \textbf{M++} will
need to calculate \textasciitilde2\textsuperscript{67} SHA-256 hash
function evaluations. Potential future optimizations of the MD5 attacks
described in \cite{ref30} have no impact on this result. This
result is achieved as follows. One can try to create a
2\textsuperscript{60}-fold multicollision for MD5 using the techniques
of Wang et al.~\cite{ref30} combined with the ideas of Joux
\cite{ref1}. This requires about 60 times more work than
creating a single collision. The length of the resulting file will be at
least 120 blocks of MD5, i.e.~7,680 bytes. Next, one searches within
this space of 2\textsuperscript{60} files for a collision for the
leftmost 120 bits of SHA-256. One expects to find such a collision by
the birthday paradox. The effort required is
\(120\cdot2^{60}\approx2^{67}\) evaluations of SHA-256.

In case a malicious user wants to forge a second preimage in an attempt
to overwrite an existing \textbf{M++} file, we can use the ideas
outlined in \cite{ref11} to show that an effort of about
\(2^{119}\) SHA-256 evaluations is required.

\subsection{Summary of naming schemes}\label{sec:5.6}\label{summary-of-naming-schemes}

The introduction of the \textbf{GM} and \textbf{M++} naming scheme
illustrates the flexibility of the Centera architecture, specifically
designed to allow for seamlessly switching amongst naming schemes.
Cryptography will continue to evolve, and due to the high importance in
the commercial and military world, hash functions will be under
continuous scrutiny. Therefore, it is important that Centera is able to
adapt its content addressing schemes to reflect the state of the art in
cryptology.

Table~\ref{tab:2} summarizes the strength of the various naming
schemes.

The second column estimates the number of files that need to be written
to a Centera system before collisions become likely.

The third column lists the effort required for an attacker to forge a
collision, i.e.~to forge two files X and \(X'\) , which result in the
same content address on a Centera. Note that this is not the same as
starting with a file X, and creating a second file \(X'\) which has the
same content address.

The fourth column lists the effort required for an attacker to forge a
second preimage. For this the attacker starts with a file X and its
content address, and the intent is to forge a second file \(X'\) which
gets the same content address on a Centera.

Note that for the \textbf{M} and \textbf{M++} naming schemes, an
attacker could attempt to create a collision or a second preimage
offline. However, for the \textbf{GM} naming scheme, all attacks need to
happen online since the attacker has no control over the 70 bits random
\textbf{G} part of the content address. This effectively makes attacks
against \textbf{GM} infeasible, since it would take millions of years to
execute these attacks.

\begin{table}[htbp]
\centering\small
\begin{tabularx}{\linewidth}{@{}l*{3}{>{\raggedright\arraybackslash}X}@{}}
\toprule
& Number of files beyond which collisions become likely & Work required to forge file collision & Work required to forge 2nd preimage of a given file $X$\\
\midrule
M & $2^{64}$ files stored & $O(1)$\textsuperscript{3} & $O(2^{108})$\\
GM & Not possible & Not possible & Not possible\\
M++ & $2^{124}$ files stored & $O(2^{67})$ & $2^{119}$\\
\bottomrule
\end{tabularx}
\par\smallskip
{\footnotesize \textsuperscript{3}For simplicity assume MD5 collisions are available from \cite{ref30}.\par}
\caption{Summary of naming schemes}\label{tab:2}
\end{table}

